% ==============================================================
\section{実験}\label{sec:experiments}
% ==============================================================

以下の実験を通じて提案手法を評価する。

\textbf{RQ1}: 空間的に局在するイベント効果を正しく帰属できるか？
\textbf{RQ2}: 1次元パイプラインに対してどの程度改善するか？
\textbf{RQ3}: イベント非対応の密集パターンを棄却できるか？
\textbf{RQ4}: 主要パラメータに対する頑健性は？
\textbf{RQ5}: 計算コストのスケーリング特性は？

\subsection{実験環境}

すべての実験はRust（release build、rayon並列化）で実装し、
Apple M3（8コア、24GB RAM）上で実行した。
デフォルトパラメータ：
$W = 200$（ウィンドウ幅）、$\theta = 3$（サポート閾値）、
$\sigma_t = W/4 = 50$、$\sigma_s = 3.0$、$B = 500$（置換回数）、$\alpha = 0.10$。

\subsection{EX1: コア帰属精度}\label{sec:ex1}

\paragraph{合成データ設計.}
$N = 10{,}000$トランザクション、10空間位置、200アイテムの語彙内データを生成する。
各トランザクションはランダムに1つの空間位置に割り当てられ、
ベースライン出現確率 $p_{\mathrm{base}} = 0.02$ でアイテムが含まれる。
3種類の信号を埋め込む：
\begin{itemize}
\item \textbf{Type A（空間局在ブースト）}: パターン $\{10, 20\}$ がイベント $E_1$ の期間中（$t \in [2000, 4000]$）、
空間位置 $\{0, \ldots, 4\}$ でブースト確率 $\beta$ で挿入される。
\item \textbf{Type B（広域ブースト）}: パターン $\{30, 40\}$ がイベント $E_2$（$t \in [5000, 7000]$）の期間中、
空間位置$\{3, \ldots, 7\}$でブーストされる。
\item \textbf{Type C（イベント非対応）}: パターン $\{50, 60\}$ がイベントとは無関係な期間
（$t \in [1000, 3000]$）・位置（$\{2, \ldots, 6\}$）でブーストされる。
\item \textbf{デコイイベント} $D_1$: パターン変動と無関係なイベント（$t \in [7500, 9000]$）。
\end{itemize}

4条件を評価する：ブースト強度 $\beta \in \{0.1, 0.2, 0.3, 0.5\}$。
各条件5シードの平均を報告する。

\paragraph{結果.}
表~\ref{tab:ex1}に結果を示す。

\begin{table}[t]
\centering
\caption{EX1: コア帰属精度（5シード平均）。
P = Precision, R = Recall, FAR = 偽帰属率。}
\label{tab:ex1}
\small
\begin{tabular}{l@{\hspace{4pt}}r@{\hspace{4pt}}r@{\hspace{4pt}}r@{\hspace{4pt}}r@{\hspace{6pt}}r@{\hspace{4pt}}r@{\hspace{4pt}}r@{\hspace{4pt}}r@{\hspace{6pt}}r}
\toprule
& \multicolumn{4}{c}{ST（提案手法）} & \multicolumn{4}{c}{1D（ベースライン）} & \\
\cmidrule(lr){2-5} \cmidrule(lr){6-9}
条件 & F1 & P & R & FAR & F1 & P & R & FAR & $\Delta$F1 \\
\midrule
$\beta{=}0.1$ & .08 & .07 & .10 & .20 & .05 & .03 & .30 & .00 & $+$.03 \\
$\beta{=}0.2$ & .47 & .50 & .50 & .00 & .09 & .05 & .40 & .00 & $+$.38 \\
$\beta{=}0.3$ & .55 & .47 & .70 & .00 & .11 & .07 & .40 & .00 & $+$.44 \\
$\beta{=}0.5$ & \textbf{.89} & \textbf{.83} & \textbf{1.0} & .00 & .09 & .06 & .20 & .00 & $+$\textbf{.80} \\
\bottomrule
\end{tabular}
\end{table}

\paragraph{分析.}
(RQ1--RQ3) $\beta \geq 0.2$ではSTパイプラインはF1で1次元パイプラインを
大幅に上回る（$\Delta F1 \geq +0.38$）。
$\beta = 0.5$ではF1 = 0.89（R = 1.00）であり、
1次元の0.09に対し+0.80の改善を達成する。
語彙内アイテムを使用した本実験では、ベースラインサポートが存在するため
1次元パイプラインは多数の偽帰属を生成しPrecisionが極めて低い。
STパイプラインは空間的近接度によりこれを効果的にフィルタリングする。

$\beta = 0.1$では両手法ともF1が低い。
STのFAR = 0.20は、信号が弱い場合に非関連パターンの
偶発的な空間一致により若干の偽帰属が発生することを示す。

\subsection{EX2: パラメータ感度}\label{sec:ex2}

EX1の$\beta = 0.3$データ（seed = 42）を用い、
主要パラメータを個別にスイープする。

\begin{table}[t]
\centering
\caption{EX2: パラメータ感度。各行は1パラメータのみ変更し他はデフォルト。}
\label{tab:ex2}
\small
\begin{tabular}{llcccc}
\toprule
パラメータ & 値 & F1 & P & R & FAR \\
\midrule
\multirow{5}{*}{$W$}
& 10  & .00 & .00 & .00 & .00 \\
& 20  & .00 & .00 & .00 & .00 \\
& 50  & .00 & .00 & .00 & .00 \\
& 100 & \textbf{1.00} & 1.00 & 1.00 & .00 \\
& 200 & .80 & .67 & 1.00 & .00 \\
\midrule
\multirow{5}{*}{$\alpha$}
& 0.01 & .00 & .00 & .00 & .00 \\
& 0.05 & .00 & .00 & .00 & .00 \\
& 0.10 & .80 & .67 & 1.00 & .00 \\
& 0.20 & .80 & .67 & 1.00 & 1.00 \\
& 0.30 & .80 & .67 & 1.00 & 1.00 \\
\midrule
\multirow{5}{*}{$B$}
& 50   & .00 & .00 & .00 & .00 \\
& 100  & .80 & .67 & 1.00 & .00 \\
& 500  & .80 & .67 & 1.00 & .00 \\
& 1000 & .80 & .67 & 1.00 & .00 \\
& 5000 & .80 & .67 & 1.00 & .00 \\
\bottomrule
\end{tabular}
\end{table}

\paragraph{分析.}
(RQ4) ウィンドウ幅 $W$ は結果に大きな影響を与える。
$W \leq 50$では信号捕捉が不十分でF1 = 0となり、
$W = 100$で最良のF1 = 1.00を達成する。
$W = 200$ではウィンドウが広すぎて変化点の時間的解像度が低下し、
1件の偽帰属が発生する。

有意水準 $\alpha$ は0.10で良好な結果が得られ、
$\alpha \geq 0.20$ではFARが上昇する。
$\alpha \leq 0.05$ではBH補正後に有意結果が消失する。
置換回数 $B$ は100以上で安定し、$B = 50$では$p$値の解像度不足により
有意結果が得られない。

\subsection{EX3: スケーラビリティ}\label{sec:ex3}

\paragraph{実験設計.}
データサイズ $N \in \{1\mathrm{K}, 5\mathrm{K}, 10\mathrm{K}, 50\mathrm{K}, 100\mathrm{K}\}$ と
イベント数 $|E| \in \{1, 3, 5, 10\}$ を変化させ、
STパイプラインと1Dベースラインの実行時間を計測する。

\paragraph{結果.}
(RQ5) 表~\ref{tab:ex3}に結果を示す。
STパイプラインの実行時間はデータサイズに対してほぼ線形にスケールし、
$N = 100{,}000$でも15秒程度で完了する（Rust release build, rayon並列化）。
1Dベースラインは一貫してSTの25--50\%の時間で完了する。
これはSTが位置ごとのサポート曲面計算と空間的近接度計算を追加で行うためである。
イベント数の増加は実行時間にほとんど影響しない（$|E| = 1$と$|E| = 10$で同等）。
候補パターン数は$N \geq 50{,}000$で飽和（19,900件）し、
パターン単位のrayon並列化により効率的に処理される。

\begin{table}[t]
\centering
\caption{EX3: スケーラビリティ。実行時間（ms）。$B = 200$。}
\label{tab:ex3}
\small
\begin{tabular}{rrccc}
\toprule
$N$ & $|E|$ & \#候補 & ST (ms) & 1D (ms) \\
\midrule
1,000  & 3  & 5     & 5      & 4     \\
5,000  & 3  & 843   & 18     & 28    \\
10,000 & 3  & 5,236 & 384    & 321   \\
50,000 & 3  & 19,900 & 6,071 & 1,870 \\
100,000 & 3 & 19,900 & 14,739 & 3,614 \\
\midrule
10,000 & 1  & 4,558 & 275    & 184   \\
10,000 & 5  & 5,539 & 678    & 292   \\
10,000 & 10 & 6,035 & 836    & 431   \\
\bottomrule
\end{tabular}
\end{table}

\subsection{考察: STパイプラインのトレードオフ}\label{sec:discussion}

実験結果から、STパイプラインと1Dパイプラインの間に明確なトレードオフが存在する。

\paragraph{精度と再現率.}
STパイプラインは語彙内信号に対して一貫して高いF1を達成し、
特にPrecisionで1Dに大きく勝る。
1次元パイプラインは空間を集約するため見かけ上のRecallは高い場合があるが、
多数の偽帰属を生成しPrecisionが極めて低い。

\paragraph{適用指針.}
位置あたりのウィンドウサイズ内トランザクション数が十分（$\geq 10$）であり、
かつイベント効果が空間的に局在する場合、STパイプラインが推奨される。
全空間で均一なイベント効果や、位置あたりの観測数が少ない場合は、
1Dパイプラインがより高いRecallを達成する可能性がある。
